\chapter{Legendre functions and spherical harmonics}\label{app:harmonics}

This appendix collects the definitions, conventions and identities used in the main text. The main reference is \citet{Varshalovich1988}, Chapters 5 and 7. Kaula's conventions are those of \citet[\S1.2]{Kaula2000}.

\paragraph{Convention used throughout (summary).} The phase convention is stated here once and kept throughout the document:
\begin{itemize}
  \item the associated Legendre functions $P_l^m$ include the Condon--Shortley phase $(-1)^m$, eq.~\eqref{eq:PCS};
  \item the spherical harmonics are $Y_l^m=N_{lm}P_l^me^{\I m\varphi}$, with no further $(-1)^m$, eq.~\eqref{eq:Ylm-def};
  \item Kaula's associated function has no phase factor, $\PK_{lm}=(-1)^mP_l^m$, eq.~\eqref{eq:PK-PCS}. Every formula of Kaula used with $P_l^m$ or $Y_l^m$ therefore carries a factor $(-1)^m$.
\end{itemize}

\section{Legendre polynomials}
The Legendre polynomials are generated by
\begin{equation}
  (1-2xt+t^2)^{-1/2}=\sum_{l\ge0}t^lP_l(x),\qquad |t|<1, \label{eq:legendre-gen}
\end{equation}
or given by Rodrigues' formula $P_l(x)=\frac{1}{2^ll!}\frac{\dd^l}{\dd x^l}(x^2-1)^l$. They obey the recursion
\[
  (l+1)P_{l+1}=(2l+1)xP_l-lP_{l-1}.
\]
For example, $P_0=1$, $P_1=x$, $P_2=\tfrac12(3x^2-1)$, $P_3=\tfrac12(5x^3-3x)$.

\section{Associated Legendre functions: two phase conventions}
\paragraph{Condon--Shortley convention (used here).} For $0\le m\le l$,
\begin{equation}
  P_l^m(x)=(-1)^m(1-x^2)^{m/2}\frac{\dd^mP_l(x)}{\dd x^m}, \label{eq:PCS}
\end{equation}
as in \citet[eq.~5.2(6)]{Varshalovich1988} and \citet{CondonShortley1935}.

\paragraph{Kaula's convention.} Without the phase factor \citep[eq.~1.28--1.30]{Kaula2000},
\begin{equation}
  \PK_{lm}(\sin\phi)=\cos^m\phi\sum_{t=0}^{\lfloor(l-m)/2\rfloor}T_{lmt}\sin^{l-m-2t}\phi,\qquad T_{lmt}=\frac{(-1)^t(2l-2t)!}{2^lt!(l-t)!(l-m-2t)!}, \label{eq:PK}
\end{equation}
where $\phi$ is the latitude, $\sin\phi=\cos\theta$. Hence
\begin{equation}
  \PK_{lm}=(-1)^mP_l^m. \label{eq:PK-PCS}
\end{equation}

The inclination functions $F_{lmp}$ and Kaula's transformation are derived with $\PK_{lm}$. Every formula of Kaula used with $P_l^m$ must therefore carry the factor $(-1)^m$ of \eqref{eq:PK-PCS}.

\paragraph{Low-degree functions} (Kaula's convention, $x=\cos\theta$, $s=\sin\theta$):
\begin{gather*}
  \PK_{20}=\tfrac12(3x^2-1),\qquad \PK_{21}=3xs,\qquad \PK_{22}=3s^2,\\
  \PK_{30}=\tfrac12(5x^3-3x),\qquad \PK_{31}=\tfrac32(5x^2-1)s,\qquad \PK_{32}=15xs^2,\qquad \PK_{33}=15s^3.
\end{gather*}
The degree-3 functions are $\PK_{3m}=s^m\mathcal P_{3m}(x)$, with $\mathcal P_{3m}=\dd^mP_3/\dd x^m$ given in~\eqref{eq:P3m}.

\paragraph{Values on the equator.} For $l-m$ even,
\[
  \PK_{lm}(0)=\frac{(-1)^{(l-m)/2}(l+m)!}{2^l\,\big(\tfrac{l-m}{2}\big)!\,\big(\tfrac{l+m}{2}\big)!},
\]
and $\PK_{lm}(0)=0$ for $l-m$ odd. For example $\PK_{20}(0)=-\tfrac12$, $\PK_{22}(0)=3$, $\PK_{31}(0)=-\tfrac32$, $\PK_{33}(0)=15$.

\section{Derivative and $1/\sin\theta$ identities}
Write $x=\cos\theta$, $s=\sin\theta$ and $P^{(m)}_l=\dd^mP_l/\dd x^m$, so that $\PK_{lm}=s^mP^{(m)}_l$. Two classical facts are used.

\paragraph{(a) Differentiated Legendre equation.} Differentiating $(1-x^2)y''-2xy'+n(n+1)y=0$ (with $y=P_n$) $k$ times gives
\[
  (1-x^2)P^{(k+2)}_n-2(k+1)xP^{(k+1)}_n+[n(n+1)-k(k+1)]P^{(k)}_n=0.
\]

\paragraph{(b) A three-term relation.} Differentiating $xP'_n-P'_{n-1}=nP_n$ $k$ times gives
\[
  xP^{(k+1)}_n+kP^{(k)}_n-P^{(k+1)}_{n-1}=nP^{(k)}_n.
\]

\paragraph{$\theta$-derivative.} Differentiating $\PK_{lm}=s^mP^{(m)}_l$ gives $\dd\PK_{lm}/\dd\theta=m\,s^{m-1}xP^{(m)}_l-s^{m+1}P^{(m+1)}_l$. Fact (a) with $n=l$, $k=m-1$ gives $s^2P^{(m+1)}_l=2mxP^{(m)}_l-(l+m)(l-m+1)P^{(m-1)}_l$. Using it for half of the second term,
\begin{equation}
  \frac{\dd\PK_{lm}}{\dd\theta}=\tfrac12\big[(l+m)(l-m+1)\PK_{l,m-1}-\PK_{l,m+1}\big]\quad(m\ge1),\qquad \frac{\dd\PK_{l0}}{\dd\theta}=-\PK_{l1}. \label{eq:dPK}
\end{equation}

\paragraph{Division by $\sin\theta$.} For $m\ge1$, divide by $s^{m-1}$; the claim
\begin{equation}
  \frac{m\PK_{lm}}{\sin\theta}=\tfrac12\big[\PK_{l+1,m+1}+(l-m+1)(l-m+2)\PK_{l+1,m-1}\big] \label{eq:mPK-sin}
\end{equation}
becomes
\[
  2mP^{(m)}_l=s^2P^{(m+1)}_{l+1}+(l-m+1)(l-m+2)P^{(m-1)}_{l+1}.
\]
Fact (a) with $n=l+1$, $k=m-1$ turns the right-hand side into $2m\big[xP^{(m)}_{l+1}-(l-m+2)P^{(m-1)}_{l+1}\big]$. Fact (b) with $n=l+1$, $k=m-1$ shows that the bracket equals $P^{(m)}_l$, which proves \eqref{eq:mPK-sin}.

The right-hand side contains no $1/\sin\theta$: the azimuthal component of the tidal force is regular at the poles.

\section{Spherical harmonics}
\begin{equation}
  Y_l^m(\theta,\varphi)=N_{lm}P_l^m(\cos\theta)e^{\I m\varphi},\qquad N_{lm}=\sqrt{\frac{2l+1}{4\pi}\frac{(l-m)!}{(l+m)!}}, \label{eq:Ylm-def}
\end{equation}
as in \citet[eq.~5.2(1)]{Varshalovich1988}. There is no further $(-1)^m$: the Condon--Shortley phase is already inside $P_l^m$. Equivalently, $Y_l^m=(-1)^mN_{lm}\PK_{lm}e^{\I m\varphi}$.

\paragraph{Properties} \citep[\S5.1]{Varshalovich1988}:
\begin{itemize}
  \item orthonormality: $\int Y_l^{m*}Y_{l'}^{m'}\dd\Omega=\delta_{ll'}\delta_{mm'}$;
  \item negative orders: $Y_l^{-m}=(-1)^mY_l^{m*}$;
  \item low degrees:
  \begin{gather*}
    Y_1^1=-\sqrt{\tfrac{3}{8\pi}}\sin\theta\,e^{\I\varphi},\qquad Y_2^0=\sqrt{\tfrac{5}{16\pi}}(3\cos^2\theta-1),\\
    Y_2^1=-\sqrt{\tfrac{15}{8\pi}}\sin\theta\cos\theta\,e^{\I\varphi},\qquad Y_2^2=\sqrt{\tfrac{15}{32\pi}}\sin^2\theta\,e^{2\I\varphi}.
  \end{gather*}
\end{itemize}

\paragraph{Addition theorem} \citep[eq.~5.17(9)]{Varshalovich1988}:
\begin{align*}
  P_l(\cos\Theta)&=\frac{4\pi}{2l+1}\sum_{m=-l}^lY_l^{m*}(\theta,\varphi)Y_l^m(\theta',\varphi')\\
  &=\sum_{m=0}^l\frac{(l-m)!}{(l+m)!}(2-\delta_{0,m})P_l^m(\cos\theta)P_l^m(\cos\theta')\cos m(\varphi-\varphi').
\end{align*}
The real form is independent of the phase convention.

\section{Recursion and differential relations}
From \citet[\S5.7--5.8]{Varshalovich1988}, with this convention:
\begin{align}
  -2m\cot\theta\,Y_l^m&=\sqrt{l(l+1)-m(m+1)}\,e^{-\I\varphi}Y_l^{m+1}+\sqrt{l(l+1)-m(m-1)}\,e^{\I\varphi}Y_l^{m-1},\\
  \partial_\varphi Y_l^m&=\I mY_l^m,\\
  \partial_\theta Y_l^m&=\tfrac12\sqrt{l(l+1)-m(m+1)}\,Y_l^{m+1}e^{-\I\varphi}-\tfrac12\sqrt{l(l+1)-m(m-1)}\,Y_l^{m-1}e^{\I\varphi},\\
  \partial_\theta Y_l^0&=\sqrt{l(l+1)}\,Y_l^1e^{-\I\varphi},\\
  \frac{m}{\sin\theta}Y_l^{m*}&=-\tfrac12\Big[c^+_{lm}Y_{l+1}^{m+1\,*}e^{\I\varphi}+c^-_{lm}Y_{l+1}^{m-1\,*}e^{-\I\varphi}\Big]\quad(m\ge1),\label{eq:Y-over-sin}
\end{align}
with $c^\pm_{lm}$ from \eqref{eq:c-constants}. The last relation is \eqref{eq:mPK-sin} multiplied by $(-1)^mN_{lm}e^{-\I m\varphi}$.

\section{Gradient of a solid harmonic}
\paragraph{Vector spherical harmonics} $\mathbf Y^{(\lambda)}_{LM}$ \citep[\S7.3, eq.~25--27]{Varshalovich1988}:
\begin{itemize}
  \item $\mathbf Y^{(1)}_{LM}=r\nabla Y_{LM}/\sqrt{L(L+1)}$, tangential;
  \item $\mathbf Y^{(0)}_{LM}=\I\vr\times\nabla Y_{LM}/\sqrt{L(L+1)}$, tangential;
  \item $\mathbf Y^{(-1)}_{LM}=\hat{\vr}Y_{LM}$, radial.
\end{itemize}

\paragraph{Gradient formulas} \citep[eq.~7.3(51)]{Varshalovich1988}:
\[
  \nabla\{fY_{LM}\}=f'\mathbf Y^{(-1)}_{LM}+\sqrt{L(L+1)}\frac fr\mathbf Y^{(1)}_{LM}
  =\sqrt{\tfrac{L}{2L+1}}\Big(\tfrac{\dd}{\dd r}+\tfrac{L+1}{r}\Big)f\,\mathbf Y^{L-1}_{LM}-\sqrt{\tfrac{L+1}{2L+1}}\Big(\tfrac{\dd}{\dd r}-\tfrac{L}{r}\Big)f\,\mathbf Y^{L+1}_{LM}.
\]
For $f=r^{-(L+1)}$ only the $\mathbf Y^{L+1}_{LM}$ term survives. Its Cartesian (cyclic) components, on $\mathbf e_{\pm1}=\mp(\mathbf e_x\pm\I\mathbf e_y)/\sqrt2$ and $\mathbf e_0=\mathbf e_z$, give the ladder relations \eqref{eq:ladder}.
